\subsection{Trees of bounded-size bags: strips, cactus graphs, and bounded treewidth}
\label{sec:classification-tree-partitions}

This is a new deduction for the classification table, recorded on
14 November 2026. It extends the binary forest construction to finite
state spaces. Its simulation input is the product-alphabet theorem in
Section~\ref{sec:resolution-product-alphabets}; it is separate from
the earlier audited revision and makes no priority claim.

\begin{theorem}[Simulation on a tree of bounded-size bags]
\label{thm:classification-bag-forest}
Partition the physical spins into known disjoint bags $B_v$, each
containing at most $b$ spins, indexed by the vertices of a forest $F$.
Suppose the classical Hamiltonian has the form
\[
 H(x)=\sum_v V_v(x_{B_v})+
       \sum_{uv\in E(F)}\Phi_{uv}(x_{B_u},x_{B_v}),\qquad
 |\Phi_{uv}(s,t)|\le B.
 \tag{BF1}
\]
The finite intrabag potentials $V_v$ are unrestricted. For fixed $b$
and fixed $\beta B<\infty$, the Gibbs family has full acceptance
simulation uniformly in the number of bags, their degrees, and all
intrabag potentials. The result is hereditary under pinnings.

More explicitly, put $\rho=1-e^{-4\beta B}$, and take $L=1$ if
$\rho=0$; otherwise set
\[
 L=\max\left\{1,\left\lceil
       \frac{2\log(8T/\epsilon)}{\log(1/\rho)}\right\rceil\right\}.
\]
For $T\ge1$ and $0<\epsilon<1$, a valid physical-spin query bound is
\[
 O\!\left(b((L+1)T)^6\epsilon^{-15}\right).
 \tag{BF2}
\]
There is no assertion of efficient computation or a uniform
zero-temperature bound.
\end{theorem}

\begin{proof}
Root the forest and regard each bag state as a single symbol in
$\Omega_v=\{-1,1\}^{B_v}$. Summing over descendant bags gives
\[
 \mu(x)=\prod_{r\text{ root}}\pi_r(x_{B_r})
             \prod_{e=(u,v)}K_e(x_{B_u},x_{B_v}),\qquad
 K_e(s,t)=\frac{w_v(t)e^{-\beta\Phi_e(s,t)}}{Z_e(s)},
 \tag{BF3}
\]
where $w_v(t)>0$ includes all descendant partition functions and
$V_v$. For any two parent states $s,s'$, (BF1) gives
\[
 e^{-4\beta B}\le\frac{K_e(s,t)}{K_e(s',t)}
                         \le e^{4\beta B}.
 \tag{BF4}
\]
Indeed the ratio of unnormalized weights is in
$[e^{-2\beta B},e^{2\beta B}]$, and the ratio of normalizers is
in the same interval. The factors $w_v(t)$ cancel, which is why
arbitrary fields and intrabag energies cause no loss.

Define the common part of the rows by
\[
 m_e(t)=\min_sK_e(s,t),\quad q_e=\sum_t m_e(t),\quad
 R_e(s,t)=\frac{K_e(s,t)-m_e(t)}{1-q_e}.
\]
When $q_e=1$, simply use a parent-independent update. Otherwise,
independently sample an edge source symbol consisting of
\[
 C_e\sim\operatorname{Bernoulli}(1-q_e),\quad
 Y_e\sim m_e/q_e,\quad
 (A_{e,s})_{s\in\Omega_u},\qquad A_{e,s}\sim R_e(s,\cdot),
 \tag{BF5}
\]
with all entries independent. Its output on parent state $s$ is
$Y_e$ if $C_e=0$ and $A_{e,s}$ if $C_e=1$.
Choose root states independently with laws $\pi_r$.
All these sources form a product law on finite alphabets. A source
query returns the whole symbol, including its table; its potentially
large computational representation is immaterial to the query model.
This constructs exactly (BF3).

Conditioned on the entire physical configuration $X=x$, source
symbols remain independent across edges and roots: each edge has
just its own constraint that the output on $x_{B_u}$ equals
$x_{B_v}$. Its posterior requires only those two bag states. In
particular,
\[
 \Pr(C_e=1\mid X=x)
 =1-\frac{m_e(x_{B_v})}{K_e(x_{B_u},x_{B_v})}
 \le1-e^{-4\beta B}=\rho<1.
 \tag{BF6}
\]
This posterior is explicit. On $C_e=0$, fix $Y_e=x_{B_v}$ and
leave the residual table at its prior law. On $C_e=1$, fix the
selected table entry $A_{e,x_{B_u}}=x_{B_v}$, leaving $Y_e$ and
all other table entries at their prior laws. Select these two cases
with their posterior probabilities. Impossible cases have zero
probability and need not be sampled. A root posterior is its observed
bag state. Thus each adaptive source request costs at most $2b$
physical queries; cache its answer.

To evaluate an output bag, follow its unique ancestor path for at
most $L$ edges. Stop at the first reset $C_e=0$ or at a root, then
apply the already read tables downwards. If neither occurs, return
a fixed bag state. This is a deterministic map on the whole product
alphabet and uses at most $L+1$ source queries. It yields a coherent
physical-bit oracle with at most $2(L+1)$ source queries by computing
and uncomputing the bag state. Given $X=x$, it can err only if all
$L$ edges on the inspected path have $C_e=1$. The path is fixed by
the forest, and posterior edge independence and (BF6) give
\[
 \Pr[F_L(Z)_i\ne x_i\mid X=x]\le\rho^L.
\]
Apply product-alphabet simulation at mean error $\epsilon/2$ to
$Q\circ F_L$, and supply its source answers with the local posterior
above. The extra hybrid error is at most
$4T\rho^{L/2}\le\epsilon/2$. Multiplying the source query bound
by $2b$ gives (BF2). Fix all auxiliary randomness by averaging.

After pinning, restrict each bag alphabet to the consistent states,
retaining singleton alphabets for fully pinned bags. The same
forest factorization and cross-potential bound hold, so the argument
is uniform under pinnings.
\end{proof}

\begin{corollary}[Classical Ising consequences]
\label{cor:classification-easy-graphs}
For pairwise Ising interactions with $|J_{ij}|\le J_*$, arbitrary
signs, and arbitrary finite fields, each of the following classes
has full AA at every fixed $\beta J_*<\infty$:
\begin{enumerate}
\item Fixed-width strips and ladders. On a nearest-neighbor strip
of width $w$, use columns as bags: $b=w$, $B\le wJ_*$.
For a periodic longitudinal direction, fold the cycle of columns
into a path of paired columns: $b\le2w$, $B\le2wJ_*$.
\item One-dimensional finite-range models with fixed range $R$.
For an open chain, consecutive groups of $R$ spins give $b=R$ and
$B\le R^2J_*$. Periodic chains may again be folded with a constant
factor increase. This also covers bounded-strength higher-order
terms of fixed spatial range, provided the total absolute strength
crossing each bag boundary is uniformly bounded.
\item Cactus graphs, with no degree bound and arbitrarily many
cycles of arbitrary lengths: $b\le2$, $B\le2J_*$.
\item Graphs with uniformly bounded treewidth and maximum degree.
\end{enumerate}
The bounds are uniform after pinnings. The strip width, range,
treewidth, and maximum degree in their respective assertions must
be independent of system size.
\end{corollary}

\begin{proof}
For a pairwise Hamiltonian on any tree-partition, use
$\Phi_{uv}(s,t)=-\sum_{i\in B_u,j\in B_v}J_{ij}s_it_j$;
its norm is bounded by the sum of crossing absolute couplings.
This proves the first two assertions using the indicated partitions.

For the third, a cactus has only edges and simple cycles as its
biconnected blocks. Root its block-cut structure. A bridge creates
a singleton child bag. For a cycle attached at an already assigned
vertex $v_0$, write its other vertices in cyclic order
$v_1,\ldots,v_{\ell-1}$ and create successive bags
$\{v_1,v_{\ell-1}\},\{v_2,v_{\ell-2}\},\ldots$, ending with
a singleton when needed. The first new bag attaches to the bag
containing $v_0$; subsequent bags form a path. Each interface carries
at most two original edges. Recursively attach all remaining blocks
at the bags containing their articulation vertices. Bags are disjoint,
have size at most two, and their quotient is a forest. This proves
the claimed constants, including for many disconnected components.

For the fourth, the graph-theoretic input is that bounded treewidth
and bounded maximum degree imply bounded tree-partition width;
see \href{https://arxiv.org/abs/math/0602507}{Wood, \emph{On
Tree-Partition-Width}, Theorem 1}. In particular, a width
$b=O((\operatorname{tw}(G)+1)\Delta(G))$ partition exists. Every
interface has at most $b\Delta(G)$ edges, so take
$B=b\Delta(G)J_*$. Apply Theorem~\ref{thm:classification-bag-forest}.
This imports a graph partition theorem, not a quantum simulation
theorem. Computation of the partition is free in our model.
\end{proof}

One may also query a fixed number $r$ of exceptional spins first,
if their deletion leaves one of these classes. Interactions with
queried spins become fields, giving an additional cost of $r$,
just as in the feedback-set corollary. Couplings incident only
through this exceptional set need not obey the remaining-edge bound.

These conclusions do not classify unbounded-width two-dimensional
grids, all bounded-treewidth graphs without degree control, arbitrary
dimension-dependent cooling schedules, or general noncommuting
quantum Hamiltonians on the same graphs. Locally diagonal classical
Hamiltonians are included when measured in their diagonal product
basis. A general quantum measurement law need not have (BF3).

The finite checks in \path{../checks/tree_partition_checks.py}
verify kernel reconstruction, posterior factorization on a small
bag tree, and the cactus partition construction. They supplement,
rather than replace, the argument above.
